% A complete unintegrated refinement; the parent assigns its statement number.
\subsection{The two dihedral images of every ordinary point-three frame}

\begin{statement}{A.D8}{Proved}
{The actual finite images of all twenty-four class-two lifts}
Let \(\mathcal G_{3,24}\) be the set of ordinary completed
unit-equation pairs from A.38 satisfying
\[
                A=3U^{24},\qquad B=-2V^{24}
                \quad(U,V\in\widehat{\Q^\times}).
\tag{D8.1}
\]
Use the actual arithmetic dihedral quotient of A.210.4.
For each pair in \(\mathcal G_{3,24}\), the twenty-four
geometric classes of integral global class-two lifts have
exactly two geometric section classes in this quotient.
They are the image of the point \(3\) and its central
twist by the quadratic character \(\kappa_2(2)\).
Each image occurs twelve times.

The distinguished lift selected by the two local coordinates
in A.40.4 has the point-three dihedral image.  Consequently
every locally point-matching class-two lift of such a pair
has this same finite image at every place.  This describes
its necessary finite image; existence of a locally matching
class-two lift is still tested by the actual localization
map of A.40.5.
\end{statement}

\begin{proof}
\emph{Ordinariness supplies the exact inertia parity.}
Use the tangential splitting, Hall coordinates, and notation
of A.210.3--6.  Thus \(C=[x,y]\), the quotient sends
\(x\) to \(s\), \(y\) to \(r\), and \(C\) to \(r^2\), and
at an odd physical prime \(q\) the actual central inertia is
\[
                  c=\frac{B_qm-A_qn}{q-1}-\frac{mn}{2}.
\tag{D8.2}
\]
At a prime \(q>3\), ordinary local realization identifies
\(m,n\) with the integer valuations of \(t_q,1-t_q\).
Equation (D8.1) puts both integers in
\(24\widehat\Z\cap\Z=24\Z\).  Their possible nonzero
patterns are \((E,0),(0,E),(-E,-E)\), with
\(E>0\) divisible by twenty-four.

For \((m,n)=(E,0)\), the actual local value
\(1-t_q\) lies in \(1+q\Z_q\).  Multiplication by two
is an automorphism of this principal-unit group, so its
full coefficient-two Kummer class is zero.  For the
chosen cocycle there is therefore \(v\in\Z_2\) with
\(B_q=(q-1)v\), and (D8.2) is \(c=Ev\).
For \((m,n)=(0,E)\), the same argument gives
\(A_q=(q-1)v\) and \(c=-Ev\).

For \((m,n)=(-E,-E)\), the literal identity
\[
            \frac{t_q}{1-t_q}
                   =-\bigl(1-t_q^{-1}\bigr)^{-1}
\tag{D8.3}
\]
and the same principal-unit argument give
\(\kappa_2(t_q)-\kappa_2(1-t_q)=\kappa_2(-1)\).
Choose the compatible roots of \(-1\) in the cyclotomic
tower; its actual cocycle is
\(\kappa_2(-1)(g)=(\chi_2(g)-1)/2\).
Allowing the integral coboundary for the already chosen
representatives gives
\[
 A_q-B_q=(q-1)(u+1/2),\qquad
 c=Eu-\frac{E(E-1)}2,\qquad u\in\Z_2.
\tag{D8.4}
\]
Each of the three formulas is even as soon as \(4\mid E\).
For \(m=n=0\), (D8.2) gives \(c=0\).
Thus in every case \(m\) is even and \(n+2c\) is
divisible by four.  The entire dihedral inertia is trivial
at every \(q>3\).  This argument uses the actual ordinary
coordinates and needs no additional divisibility of
\(E/(q-1)\).

\emph{The finite conjugation map.}
Choose representatives using the completed roots in
(D8.1).  After geometric conjugation the two abelian
cocycles of any global lift agree modulo two literally
with those of the point-three section.  Their dihedral
difference is a central quadratic character, and the
preceding inertia computation, together with the same
zero-valuation computation for the point \(3\), shows
that it is represented by
\[
                   d\in\{1,-1,2,-2,3,-3,6,-6\}.
\tag{D8.5}
\]

We calculate every geometric conjugator in the actual
arithmetic quotient.  Write the point-three cocycle as
\(s^{a(g)}r^{b(g)}\), retaining its actual exponent
\(b(g)\) modulo four.  Its reductions are
\[
 a(g)\bmod2=\kappa_2(3)(g),\qquad
 b(g)\bmod2=\kappa_2(-2)(g).
\]
The tangential arithmetic action is
\(g(s)=s\), \(g(r)=r^{\chi_2(g)}\).
Conjugation by \(s\) changes this section by
\((r^2)^{b(g)}\).  Conjugation by \(r\) changes it by
\((r^2)^{a(g)+(\chi_2(g)-1)/2}\).
The full central character changes are therefore
\[
 \begin{aligned}
       \delta_s&=\kappa_2(-2),\\
       \delta_r&=\kappa_2(3)+\kappa_2(-1)
                                      =\kappa_2(-3),\\
       \delta_{r^2}&=0.
 \end{aligned}
\tag{D8.6}
\]
These formulas may also be checked directly in
\(s^2=r^4=1,\ srs=r^{-1}\); they use both the
point-three cocycle and the arithmetic action.
The central commutator has trivial arithmetic action,
so conjugation gives an additive character map and
these three formulas account for all eight geometric
elements.  The same map applies after any central twist.

The characters \(-2,-3\) are independent rational
squareclasses.  Thus two of the supported twists
(D8.5) are geometrically conjugate exactly when their
ratio lies in
\[
                           \{1,-2,-3,6\}.
\tag{D8.7}
\]
This subgroup has index two in (D8.5), and \(2\)
represents its nonidentity coset.  Hence the two
possible geometric classes are precisely the asserted
point-three class and its twist by \(2\).

\emph{The actual global correction and the selected class.}
By A.40 the twenty-four global lift classes form a
torsor under \(\Z/24\), with representative cocycle
\(\eta_2=(\chi^2-1)/24\).  Its reduction in the
dihedral centre is
\[
             \eta_2\bmod2=\kappa_2(2).
\tag{D8.8}
\]
Indeed its values for \(\chi\bmod8=1,3,5,7\) are
respectively \(0,1,1,0\); these are the actual
automorphism signs of
\(\sqrt2=\zeta_8+\zeta_8^{-1}\).
Consequently even and odd corrections give the two
distinct classes in (D8.7), twelve times each.

To identify the distinguished lift, work at the
physical prime three.  Ordinariness and (D8.1) give
\[
            v_3(t_3)\equiv1,\qquad
            v_3(1-t_3)\equiv0\pmod {24}.
\tag{D8.9}
\]
The first valuation is not zero.  If it were negative,
the two valuations would be equal, contradicting
their displayed residues.  Thus \(v_3(t_3)>0\).
The actual punctured-disc comparison of A.49 compares
the point sections of \(t_3\) and \(3\) with the
cusp-zero sections of those tangent parameters.
In this dihedral quotient that cusp inertia has
order two; (D8.1) makes the two tangent Kummer
characters equal modulo two.  Their local dihedral
section classes therefore agree.  No equality of
the full pro-two tangent parameters is required.

Finally the two global classes in (D8.7) remain
distinct after restriction to \(\Q_3\).  There
\(-2\) is a square, \(-3\) has odd valuation and
defines a ramified quadratic character, while
\(2\) is a nonsquare unit and defines the
nontrivial unramified quadratic character.
The exact local conjugation map (D8.6) therefore
cannot identify the point-three class with its
twist by \(2\).  The distinguished lift of A.40
matches the actual coefficient-two class-two point
section at three, so its global dihedral class
must be the point-three class.  The same applies
to every lift matching the ordinary local points
at all finite places.
\end{proof}

The exact certificate
\path{code/point3_dihedral_frame_certificate.py}
checks the Hall substitutions, every dihedral
conjugator and odd arithmetic action, the rational
quadratic-character quotient, its local restriction
at three, and the parity of all twenty-four
central corrections.  The ordinary-prime
argument and the punctured-disc map are proved
above and are not inferred from finite sampling.
